%% kapa: exam=1 round=32 year=2021 period=2 form=A subject=accounting subject_ko=회계학 seq=27 number=67
%% source: Q-Net 감정평가사 기출문제 (cst00309) artlSeq=5208397 file=extracted/1008/5208397/A형/2021년도 제32회 감정평가사 1차 2교시 A형.pdf
%% conversion: kordoc 4.24.51 -> tools/gen1.py (auto); review=proofread-sonnet; flags=none
\begin{question}{67}
(주)감평은 재고자산을 20×1년 말까지 평균법을 적용해 오다가 20×2년 초 선입선출법으로 회계정책을 변경하였다. 다음은 20×1년 말과 20×2년 말의 평가방법별 재고자산 금액이다.

{\small\centering\begin{tabular}{|c|c|c|c|}\hline
\multicolumn{2}{|c|}{구분} & 20×1년 말 & 20×2년 말 \\ \hline
\multirow{2}{*}{재고자산 금액} & 평균법 & ￦2,800 & ￦2,200 \\ \cline{2-4}
 & 선입선출법 & 2,500 & 2,800 \\ \hline
\end{tabular}\par}
\vspace{4pt}
평균법을 적용한 20×2년 당기순이익이 ￦2,000일 때, 변경 후 20×2년 당기순이익은? (단, 동 회계정책 변경은 한국채택국제회계기준에서 제시하는 조건을 충족하는 것이며, 선입선출법으로의 회계정책 변경에 대한 소급효과를 모두 결정할 수 있다고 가정한다.)
\begin{choices}
  \item ￦1,400
  \item ￦2,000
  \item ￦2,300
  \item ￦2,600
  \item ￦2,900
\end{choices}
\end{question}
\answer{5}
